SQL:2016 introduced the \texttt{MATCH\_RECOGNIZE} clause to search ordered
rows for user-defined patterns.  It is useful when the information is a
sequence, not one row.  It expresses V-, W-, and U-shaped patterns
(one dip, two dips, a long dip) and repeated behavior, as in fraud
detection, IoT monitoring, and clickstreams.
Oracle and Trino provide it, but Pandas does not, so analysts fall back on
shifted columns, loops, and state handling that are harder to verify.

The standard defines row-pattern recognition as three independent optional
features: R010 places the clause in the \texttt{FROM} clause, so a pattern
reads like a table, and the other two cover windows and aggregates.  A
vendor may implement any or none, so
support differs widely.

This thesis addresses that gap with an in-process engine for a subset of
R010 that it defines, implements, and tests.  The engine parses a query,
compiles its row pattern into finite automata, and executes it over ordered
DataFrame partitions.  The evaluation covers correctness, feature
support, usability, performance, and scalability against Trino and Oracle.
All 709 regression tests passed.

The main benchmark ran five patterns over six Amazon UK Products~2023 sizes,
100\,K to 2.22\,million rows (30 pattern--size cells per system), and all
three agreed on
every result.  Under one core and 32\,GB, mean time was
0.17\,s for the engine, 1.04\,s for Oracle, and 1.98\,s for Trino.  A scalability study on Amazon Reviews~2023 ran the same five patterns over
seven sizes to 227.9\,million rows (35 cells per system) under one core and
58\,GB; the engine and Oracle completed all 35, Trino 33.
Engine time stayed close to linear, while query-time memory grew with input size.
Row-pattern recognition can therefore connect declarative SQL with in-process
Pandas analysis while keeping sequential queries readable.

\noindent\textbf{Keywords:} \texttt{MATCH\_RECOGNIZE}; Row Pattern
Recognition; Pandas; SQL:2016; Finite Automata; Pattern Matching.
